\section*{Chapitre \ref{ch:iv2:macro-market-knowledge-and-pitches} --- Macro, connaissance des marchés et pitches}
\markboth{\protect\omTocMark{}{Réponses~: Macro et pitches}}{\protect\omTocMark{}{Réponses~: Macro et pitches}}

\begin{solution}{iq:iv2:macro-market-knowledge-and-pitches:1}
L'obligation perd environ 7~\%, et l'action environ 20~\%. Avec une croissance en hausse d'un demi-point, l'action baisse d'environ 11~\%.

Le cours de l'action par dollar de dividende passe de $1/0.04 = 25$ à $1/0.05 = 20$. Avec $g = 4.5~\%$, il vaut $1/0.045 = 22.2$.

Une action à duration longue peut donc bouger davantage qu'une obligation, et la cause de la variation des taux détermine une bonne part du résultat.

\iqwrong{Affirmer que la hausse des taux signifie toujours la baisse des actions. Si les taux montent sur une croissance plus forte, les bénéfices peuvent plus que compenser cet effet.}
\end{solution}

\begin{solution}{iq:iv2:macro-market-knowledge-and-pitches:2}
Il vaut 2.5~\%. C'est le taux d'inflation auquel les deux obligations rapportent la même chose. Il contient aussi une prime de risque d'inflation et un écart de liquidité.

$4.2 - 1.7 = 2.5$. Si l'inflation moyenne dépasse ce niveau, l'obligation indexée gagne.

Ce n'est donc pas une simple prévision. Une prime de risque le ferait monter au-dessus de l'inflation anticipée, et la moindre liquidité de l'obligation indexée le fait baisser.

\iqwrong{Le lire comme la prévision d'inflation du marché sans correction. Il en est proche, mais il peut varier pour des raisons techniques en période de tension.}
\end{solution}

\begin{solution}{iq:iv2:macro-market-knowledge-and-pitches:3}
Les taux courts montent, les taux longs montent en général moins, la devise tend à s'apprécier, les actions tendent à baisser, et l'or a un comportement mitigé. Tout cela est relatif à ce qui était intégré dans les prix.

La chaîne~: plus d'inflation signifie une politique monétaire plus restrictive que prévu. Les rendements à court terme montent donc le plus, et la courbe tend à s'aplatir. Des taux plus élevés attirent les capitaux, la devise monte donc. Les actions baissent avec le taux d'actualisation plus élevé, et les secteurs sensibles aux taux baissent le plus. L'or ne verse aucun coupon~: il baisse avec les rendements réels plus élevés, même s'il peut profiter des craintes d'inflation.

Pour construire sa propre réponse, il faut toujours préciser quel canal on vise~: politique anticipée, compensation d'inflation ou taux réels. Puis dire ce qui pourrait l'invalider~: une banque centrale perçue comme en retard sur la courbe ferait plutôt monter les rendements longs, et la pente de la courbe s'accentuerait.

\iqwrong{Donner une liste de sens de variation sans canal. Il faut dire pourquoi. Quelqu'un qui raisonne à partir d'un mécanisme peut s'adapter quand les données changent, alors qu'une liste apprise par cœur ne le peut pas.}
\end{solution}

\begin{solution}{iq:iv2:macro-market-knowledge-and-pitches:4}
Une probabilité de 52~\% pour une baisse de taux était intégrée dans les cours. Un maintien fait monter les taux courts d'environ 13 points de base. Les actions tendent à baisser sur cette surprise. Ce qui fait bouger les cours, c'est l'écart avec ce qui était intégré.

$(4.33 - 4.20)/0.25 = 0.52$. Après un maintien, le taux attendu monte donc de 4.20 à 4.33.

Si un maintien avait été intégré à 100~\%, rien ne se passerait.

\iqwrong{Dire «~pas de changement, donc pas de réaction~». Une absence de changement du taux peut représenter un fort changement des anticipations.}
\end{solution}

\begin{solution}{iq:iv2:macro-market-knowledge-and-pitches:5}
Le chiffre d'affaires publié recule de 8~\%. Ensuite je demanderais des précisions sur les coûts~: si leur part en euros est supérieure à celle des ventes, le bénéfice baisse moins. Je demanderais aussi comment l'entreprise se couvre, et si des prix plus bas en dollars leur font gagner des ventes.

$0.8 \times 10~\% = 8~\%$. Les 20~\% restants ne bougent pas.

\iqwrong{Supposer que le bénéfice baisse aussi de 8~\%. Cela n'est vrai que si les coûts se répartissent en euros dans la même proportion que les ventes. Avec davantage de coûts en euros, le bénéfice baisse moins, et une base de coûts en dollars est bien plus touchée.}
\end{solution}

\begin{solution}{iq:iv2:macro-market-knowledge-and-pitches:6}
Il faut préparer une courte liste avant chaque entretien, avec la date, et le dire quand on la donne. Si l'on ne sait pas, le dire, donner la fourchette dont on est sûr, et indiquer comment on la vérifierait.

\textbf{Pour se préparer.} Il faut rédiger une liste d'une page et la mettre à jour le matin de l'entretien. Taux~: taux directeur, rendements des emprunts d'État à deux ans et à dix ans. Actions~: indice principal et son mouvement récent. Volatilité~: volatilité implicite de l'indice principal. Devises~: les principales paires de devises en dollar. Matières premières~: pétrole brut, gaz naturel, or. Crédit~: un spread investment grade et un spread high yield. Il faut noter le mouvement marquant du dernier mois pour chacun, et une raison de ce mouvement.

\textbf{Ce qu'il faut dire.} «~À ce matin, environ X. Il a varié d'environ Y ce mois-ci à cause de Z.~» Si l'on n'est pas sûr~: «~Je dirais dans une fourchette de A à B, mais il faudrait vérifier.~» Une fourchette large, honnêtement donnée, vaut mieux qu'un chiffre faux donné avec assurance.

\textbf{Pour construire sa propre liste.} Il faut garder une source pour chaque ligne, comme la page d'une banque centrale ou les données d'une bourse, afin de pouvoir indiquer d'où elle provient.

\begin{dated}{2026-11}{Dates à connaître}\label{dat:iv2:macro-market-knowledge-and-pitches:cal}
Aux États-Unis, la Fed tient huit réunions de politique monétaire par an, et les rapports mensuels sur l'emploi et les prix à la consommation sortent à 8~h~30, heure de l'Est. Les dates changent~: une interruption de financement de l'administration fédérale a annulé un rapport sur l'emploi en 2025. Il faut consulter le calendrier en vigueur.
\end{dated}

\iqwrong{Bluffer sur un niveau. Les recruteurs peuvent le vérifier en un instant, et un chiffre faux donné avec assurance coûte plus cher que de dire qu'on ne sait pas.}
\end{solution}

\begin{solution}{iq:iv2:macro-market-knowledge-and-pitches:7}
Cela représente environ 238~millions de dollars, soit environ 24~\% du bénéfice. Si la moitié est couverte, c'est environ 12~\%.

$\$10/42 = \$0.238$ par gallon, sur 1 milliard de gallons. La moitié couverte est protégée, donc c'est 119~millions de dollars.

\iqwrong{Oublier les tarifs. Si les concurrents subissent le même coût, une partie se répercute sur les passagers. C'est pourquoi le chiffre du premier ordre est un plafond, mais seulement si la demande tient.}
\end{solution}

\begin{solution}{iq:iv2:macro-market-knowledge-and-pitches:8}
Cela implique environ 3.875~\% sur les huit ans. Une courbe inversée traduit une anticipation de taux courts plus bas plus tard, ou une prime de terme plus faible. Elle ne dit pas qu'une récession est certaine.

$((1.04)^{10}/(1.045)^2)^{1/8} - 1 = 3.875~\%$. La version rapide~: $(10 \times 4.0 - 2 \times 4.5)/8 = 3.875~\%$.

Les inversions passées ont souvent précédé des récessions, sur un petit nombre d'épisodes, mais le moment où le signal se produit est flou, et les variations de la prime de terme peuvent le fausser.

\iqwrong{Prendre la corrélation avec les récessions passées pour une loi causale. Il s'agit d'une affirmation sur les trajectoires de politique monétaire intégrées dans les prix, tirée de quelques cas seulement.}
\end{solution}

\begin{solution}{iq:iv2:macro-market-knowledge-and-pitches:9}
L'action baisse d'environ 5~\%~: $+5~\% - 0.5 \times 20 = -5~\%$. Ils montent ensemble lorsque la hausse des rendements vient d'une croissance plus forte, et que les bénéfices progressent plus que ne pèse le taux d'actualisation.

Une duration de 20 signifie qu'un point de taux d'actualisation représente 20~\% de la valeur. Donc 50 points de base correspondent à 10~\%.

Dans une période où l'inflation est bien ancrée, des rendements plus élevés signifient souvent une meilleure croissance, la corrélation est donc positive. Quand l'inflation est la préoccupation, des rendements plus élevés signifient une politique plus restrictive, et la corrélation devient négative.

\iqwrong{Considérer que la relation entre les deux est fixe. Le signe dépend de ce qui pousse les rendements.}
\end{solution}

\begin{solution}{iq:iv2:macro-market-knowledge-and-pitches:10}
Elle donne 6~\%, et 7.5~\% avec une inflation d'un point plus élevée.

$2 + 3 + 0.5 \times 1 + 0.5 \times 1 = 6$. Avec $\pi = 4$~: $2 + 4 + 1 + 0.5 = 7.5$. Donc une hausse d'un point de l'inflation fait monter le taux directeur de 1.5 points, ce qui fait monter les taux réels.

Elle ne dit pas ce qui est intégré dans les prix, et l'écart de production et $r^*$ ne sont pas observables. Une règle est un indice de référence pour mesurer à quelle distance se trouve la politique de son comportement passé, et elle ne tranche rien.

\iqwrong{Affirmer que la règle prouve que la politique est erronée. Des estimations différentes de $r^*$ et de l'écart peuvent déplacer le résultat d'un point ou plus.}
\end{solution}

\begin{solution}{iq:iv2:macro-market-knowledge-and-pitches:11}
Elle vaut environ 9.2~\% pour $-0.3$, et environ 10.8~\% pour $+0.3$.

La variance vaut $0.36 \times 0.0256 + 0.16 \times 0.0049 + 2 \times 0.24 \times \rho \times 0.0112 = 0.01 + 0.00538\rho$.

Le signe dépend de ce qui pilote le marché. Si les chocs de croissance dominent, les obligations montent quand les actions baissent, donc la corrélation est négative. Si les chocs d'inflation et de politique monétaire dominent, la hausse des taux pèse sur les deux, donc la corrélation est positive.

\iqwrong{Supposer que les obligations sont toujours une couverture. Elles l'ont été pendant une longue période à faible inflation, mais ce n'est pas une loi.}
\end{solution}

\begin{solution}{iq:iv2:macro-market-knowledge-and-pitches:12}
Cela représente environ 5~points de base, soit une perte de 50\,000~\$. Les rendements pourraient ne pas bouger si le chiffre était déjà intégré dans les cours grâce à d'autres données, si les révisions et les salaires sont faibles, ou si le taux de chômage dit le contraire.

$150\,000/30\,000 = 5$. Et $5 \times \$10\,000$.

Les payrolls font l'objet de révisions importantes, et l'enquête auprès des ménages peut donner un autre résultat. Une bonne surprise peut donc être annulée au sein même de la publication. Le bêta est une moyenne sur de nombreux régimes, et il change lorsque les priorités de la Fed évoluent.

\iqwrong{Trader sur le seul titre. Lire la publication~: révisions, salaires, chômage et taux de participation. Et vérifier ce qui était déjà intégré après les données précédentes.}
\end{solution}

\begin{solution}{iq:iv2:macro-market-knowledge-and-pitches:13}
Le taux à deux ans baisse d'environ 15~points de base, et le taux à dix ans d'environ 5, donc la courbe se pentifie de 10, et le steepener rapporte environ 100\,000~\$. Le reste dépend de la raison pour laquelle la banque a été si audacieuse.

La surprise est de 25. Donc $0.6 \times 25 = 15$, et $0.2 \times 25 = 5$. La position longue à deux ans gagne $15 \times \$10\,000$, et la position courte à dix ans perd $5 \times \$10\,000$.

\textbf{Le reste.} La devise a tendance à baisser. Les actions peuvent monter grâce à une politique monétaire plus accommodante. Elles peuvent aussi baisser si le marché lit l'ampleur de l'action comme une panique~: «~qu'est-ce qu'ils savent que nous ne savons pas~?~» C'est tout l'intérêt d'une question sur un événement~: énoncer les deux lectures, et dire ce qui permettrait de les départager.

\iqwrong{Prendre les 50 en entier comme choc. Seule la surprise de 25 constitue une information.}
\end{solution}

\begin{solution}{iq:iv2:macro-market-knowledge-and-pitches:14}
Le dollar peut baisser d'environ 3.8~\% avant qu'une perte n'apparaisse. Le ratio de Sharpe est d'environ 0.4. Dans un débouclage, le carry ne protège pas, et une forte baisse touche de nombreux détenteurs à la fois.

$1.045 \times x/1.005 = 1$ donne $x = 0.9617$. Cela fait $-3.83~\%$. Le carry de 4~\% pour une volatilité de 10~\% donne un ratio de Sharpe de 0.4.

Une baisse d'un écart-type du dollar fait perdre environ 6~\%. Comme beaucoup de détenteurs ont la même position avec effet de levier, une baisse force les autres à vendre, et cela rend la queue de la distribution des rendements plus épaisse~: de petits gains pendant des années, puis une perte soudaine et importante.

\iqwrong{Prendre le forward pour une prévision, ou dire que le carry est gratuit. Un ratio de Sharpe de 0.4 est un rendement modeste pour un skew très négatif.}
\end{solution}

\begin{solution}{iq:iv2:macro-market-knowledge-and-pitches:15}
Elle modifie l'offre d'obligations longues que les investisseurs privés doivent détenir, ce qui agit sur leur prime de risque, et elle signale la trajectoire des taux courts. Elle est difficile à mesurer, car les deux effets se produisent en même temps et les annonces sont anticipées.

Il existe deux canaux. \textbf{Effet d'équilibre de portefeuille~:} acheter de la duration laisse moins d'obligations aux détenteurs privés, qui acceptent donc une prime de terme moindre. \textbf{Effet de signal~:} une banque centrale qui achète des obligations est aussi peu susceptible de relever ses taux à court terme. Les deux effets font baisser les rendements longs, et l'inverse vaut lorsqu'elle vend.

Pour construire la réponse~: dire ce qui relève de l'effet de stock (le montant détenu au total) et ce qui relève de l'effet de flux (le rythme). Indiquer lequel compte selon soi. Dire que les estimations viennent en général de fenêtres d'annonce, qui ne captent que les surprises et dépendent elles-mêmes du modèle.

\iqwrong{Dire «~ils impriment de la monnaie, donc l'inflation suit.~» Les réserves s'accumulent auprès des banques, et le lien avec la masse monétaire au sens large et avec les prix n'est pas mécanique.}
\end{solution}

\begin{solution}{iq:iv2:macro-market-knowledge-and-pitches:16}
Le cours recule d'environ 5.5~\%. Le marché accordait de l'importance aux prévisions, et non au trimestre écoulé.

$1.05 \times 18/20 = 0.945$.

Un chiffre supérieur aux attentes regarde vers le passé, et le multiple est un prix sur l'avenir. Ainsi, un même résultat peut donner des réactions opposées, selon ce qui est dit sur l'année suivante.

\iqwrong{Dire «~Ils ont fait mieux, donc le cours doit monter.~» Le cours intégrait déjà une attente, sur les chiffres comme sur les prévisions.}
\end{solution}

\begin{solution}{iq:iv2:macro-market-knowledge-and-pitches:17}
Adopter cet ordre~: la thèse en une ligne, la raison pour laquelle le titre est mal évalué, un chiffre pour sa valeur, le catalyseur, les risques, puis la taille de la position et la sortie. Être prêt à défendre chacun de ces points.

\textbf{Ce qui doit y figurer.}
\begin{enumerate}
\item \textbf{Thèse en une ligne.} «~Position longue sur [société fictive], car les marges vont se redresser plus vite que ne le pense le marché.~»
\item \textbf{La vision divergente.} Ce que pense le marché et pourquoi il a tort. Sans vision divergente, il n'y a pas de pitch.
\item \textbf{Une valorisation chiffrée.} Le BPA estimé face au consensus, le multiple retenu, le cours cible. Scénarios haussier, central et baissier.
\item \textbf{Un catalyseur.} Ce qui amène les autres à le voir, et à quel moment.
\item \textbf{Risques.} Nommer les deux qui pèseraient le plus, et dire ce qu'il faudrait surveiller.
\item \textbf{Taille et sortie.} Un stop ou une condition d'invalidation de la thèse, et la taille de la position.
\end{enumerate}

\textbf{Comment construire la sienne.} Choisir une société dont on peut dire quelque chose de concret, dans un secteur que l'on a étudié. Utiliser les chiffres courants issus de ses dépôts réglementaires, et les vérifier le matin même de l'entretien. Avoir prêts le multiple, le consensus et les principaux pairs, avec leur date. S'entraîner 2 minutes, puis 5. Ne pas pitcher une action parce qu'elle a monté, ou une valeur que tout le monde apprécie sans dire ce que le marché ne voit pas.

\iqwrong{Raconter l'histoire de la société sans parler du cours. «~C'est une excellente société~» n'est pas un pitch, car une excellente société à un cours élevé peut être une mauvaise position longue. Le recruteur veut savoir ce qui est déjà dans le cours, et pourquoi on s'en distingue.}
\iqfollow{Et si le recruteur vous dit que c'est déjà intégré dans le cours~?}{Montrer ce qui est déjà intégré, avec un chiffre~: le multiple implique un certain taux de croissance. Dire ensuite pourquoi son propre chiffre diffère. Si l'on n'y parvient pas, concéder, et dire ce qui ferait changer d'avis.}
\iqfollow{Et si un événement macroéconomique survient ensuite~?}{Faire passer l'événement par ses moteurs~: les taux par le multiple, la devise par la répartition du chiffre d'affaires, et les coûts d'intrants par la marge. Dire ensuite ce que l'on ferait de la position.}
\end{solution}

\begin{solution}{iq:iv2:macro-market-knowledge-and-pitches:18}
L'espérance de rendement est d'environ 9.4~\%. Le rapport récompense sur risque est d'environ 0.84 sur l'ensemble des scénarios, ou d'environ 2.5 jusqu'au stop. La position représente environ 8.3~\% du portefeuille. Un gap franchissant le stop est ce qui peut mal tourner.

\textbf{Scénarios.} Haussier~: $3.60 \times 18 = 64.80$, soit $+29.6~\%$. Central~: $3.30 \times 16 = 52.80$, soit $+5.6~\%$. Baissier~: $2.70 \times 12 = 32.40$, soit $-35.2~\%$. Le rendement espéré est $0.5 \times 29.6 + 0.3 \times 5.6 + 0.2 \times (-35.2) = 9.44~\%$.

\textbf{Rapport récompense sur risque.} Gain haussier rapporté à la perte baissière~: $29.6/35.2 = 0.84$. Jusqu'au stop, $29.6/12 = 2.47$.

\textbf{Taille.} Un stop à 6~\$ sur 50~\$ représente 12~\%. Donc $1~\%/12~\% = 8.3~\%$ du portefeuille.

\textbf{Le stop.} Les gaps sur résultats franchissent les stops. Dans le scénario baissier, un gap à la baisse jusqu'à 32.40 fait perdre 35.2~\%, soit 2.9~\% du portefeuille, environ trois fois le plan. Donc si une publication de résultats intervient avant le catalyseur, il faut soit réduire la taille, soit acheter un put.

Dire aussi quelle part des 9.4~\% dépasse ce que l'on gagnerait en détenant le marché avec le même bêta. C'est le véritable avantage.

\iqwrong{Dimensionner la position à partir du scénario haussier. C'est la perte qui fixe la taille. Un stop que l'on peut franchir d'un bond ne limite pas ce que l'on perd.}
\end{solution}

\begin{solution}{iq:iv2:macro-market-knowledge-and-pitches:19}
Les notionnels valent environ 52.6~millions de dollars et 12.5~millions de dollars. Un steepening de 25 points de base rapporte 250\,000~\$. Le carry est d'environ $-\$65\,000$ par an. La VaR est d'environ 140\,000~\$.

\textbf{Notionnels.} $10\,000/(1.9 \times 0.0001) = 52.6$ millions de dollars, et $10\,000/(8 \times 0.0001) = 12.5$ millions de dollars.

\textbf{Gain.} Le spread varie de 25, à 10\,000~\$ le point de base~: 250\,000~\$, quel que soit le niveau parallèle. La position est neutre face à un mouvement parallèle.

\textbf{Carry.} Position longue à deux ans~: $52.6 \times (4.0 - 4.1)\% = -\$52\,600$. Position courte à dix ans~: on paie 4.2~\% et l'on perçoit le repo à 4.1~\%, donc $12.5 \times (-0.1)\% = -\$12\,500$. Au total, environ $-\$65\,000$ par an, soit environ 5\,400~\$ par mois. Sur trois mois, il faut environ 1.6~points de base de steepening pour le compenser. Le roll-down de la courbe est laissé de côté ici, et il faudrait l'ajouter lorsque la courbe présente une vraie pente à ces échéances.

\textbf{Risque.} $6 \times \$10\,000 \times 2.326 = \$139\,600$.

\textbf{Le pitch.} Pourquoi maintenant~: la trajectoire déjà intégrée dans le cours, face à sa propre vue. Stop~: dire 15 points de base contre la position, soit une perte de 150\,000~\$. Ne pas oublier que le second notionnel vaut un quart du premier, et que ce sont les DV01 qui doivent être égaux.

\iqwrong{Prendre des notionnels égaux. Cela ferait un gros pari sur le dix ans, et non sur la courbe.}
\end{solution}

\begin{solution}{iq:iv2:macro-market-knowledge-and-pitches:20}
Rendements nominaux en hausse de 1~point~: l'obligation nominale rapporte $-4~\%$, et l'obligation indexée $+4.7~\%$. Rendements réels en hausse de 1~point~: l'obligation nominale rapporte $-4~\%$, et l'obligation indexée $-3.3~\%$. Elle couvre donc l'inflation, et non les taux réels.

\textbf{Nominal.} Son portage vaut 4, et la variation de prix vaut $-8$. Soit $-4$ dans les deux cas.

\textbf{Indexée.} Son rendement vaut le rendement réel 1.7, plus l'inflation 3, moins la duration sur le rendement réel. Soit $4.7$ si le rendement réel ne bouge pas, et $4.7 - 8 = -3.3$ s'il monte de 1~point.

Ainsi, lorsque la hausse des rendements nominaux provient des points morts d'inflation, l'obligation indexée fait nettement mieux. Lorsqu'elle provient des rendements réels, elle baisse avec le reste.

\iqwrong{Affirmer qu'une obligation indexée sur l'inflation ne peut pas perdre d'argent lors d'une année d'inflation élevée. Lorsque la banque centrale réagit et que les rendements réels montent, les obligations indexées de longue duration peuvent chuter fortement.}
\end{solution}

\begin{solution}{iq:iv2:macro-market-knowledge-and-pitches:21}
Les rendements réels ont bougé, ce qui traduit une politique plus restrictive ou une meilleure croissance, avec moins d'inflation anticipée. Les obligations nominales baissent d'environ 3.75~\%, et les obligations indexées d'environ 4.8~\%. Je préférerais détenir des actions de valeur, car les actions de croissance ont la duration la plus longue.

$-7.5 \times 0.5 = -3.75$, et $-8 \times 0.6 = -4.8$.

La baisse des points morts d'inflation évoque un scénario de politique monétaire~: la crédibilité progresse, les craintes d'inflation reculent, et les taux réels montent. Ce n'est pas un pur scénario de croissance, dans lequel les points morts d'inflation monteraient en général eux aussi.

Les actions de croissance tirent l'essentiel de leur valeur de flux lointains~; elles sont donc les plus touchées par une hausse des taux réels. La valeur a des flux plus proches, et elle comprend davantage de valeurs cycliques.

\iqwrong{Ne regarder que le rendement nominal. Un mouvement de 50 points de base peut provenir de scénarios très différents, avec des gagnants différents. Il faut toujours commencer par le décomposer.}
\end{solution}

\begin{solution}{iq:iv2:macro-market-knowledge-and-pitches:22}
Une carte simple~: croissance en hausse et inflation en baisse favorise les actions. Croissance en hausse et inflation en hausse favorise les matières premières et les valeurs cycliques. Croissance en baisse et inflation en baisse favorise les obligations. Croissance en baisse et inflation en hausse pénalise presque tout, sauf les matières premières et le cash. Le danger est qu'elle masque ce qui est déjà dans les cours.

\begin{center}
\begin{tabular}{@{}>{\raggedright\arraybackslash}p{2.6cm}>{\raggedright\arraybackslash}p{6.3cm}>{\raggedright\arraybackslash}p{6.3cm}@{}}
\toprule
 & Inflation en baisse & Inflation en hausse \\
\midrule
Croissance en hausse & Les actions gagnent~; obligations correctes & Les matières premières gagnent~; actions mitigées~; obligations perdent \\
Croissance en baisse & Les obligations gagnent~; actions perdent & Cash et matières premières~; actions et obligations perdent toutes deux \\
\bottomrule
\end{tabular}
\end{center}

\textbf{Pourquoi elle est dangereuse.} Elle trie par niveau, alors que les cours réagissent aux surprises. Si une forte croissance est déjà intégrée dans les cours, un rapport simplement bon déçoit. Les actifs changent aussi~: un indice riche en énergie ne se comporte pas comme un indice riche en technologie. Et quelques décennies de données comptent peu de régimes.

\textbf{Comment construire la sienne.} Il faut confronter la carte à plusieurs années passées, avec les rendements effectifs issus d'une source publique. Être prêt à citer une année qui ne correspond pas, et la raison.

\iqwrong{Réciter le tableau comme s'il s'agissait d'une loi. Dire «~tend à~», dire que l'effet est relatif à ce qui est déjà intégré dans les cours, et dire quelle observation ferait changer d'avis.}
\end{solution}

\begin{solution}{iq:iv2:macro-market-knowledge-and-pitches:23}
Le rendement non couvert est de 4.5~\%. Couvert, il avoisine 11.5~\%. La couverture rapporte l'écart de taux autant que la baisse de l'euro.

Non couvert~: $1.10 \times 0.95 - 1 = 4.5~\%$.

Le terme à terme vaut $F/S = 1.045/1.025 = 1.0195$. Vendre des euros à terme sur la valeur de départ rapporte $1.0195 - 0.95 = 6.95~\%$. Soit $4.5 + 6.95 = 11.45~\%$.

Vendre à terme une devise à faible taux revient à la vendre avec une prime. La couverture rapporte donc le portage de 2~\%. Si l'euro ne bouge pas, on obtient 10~\% plus environ 2~\%. La couverture ne porte que sur la valeur de départ~; le gain de 10~\% reste donc exposé, ce qui explique le faible résidu ici.

\iqwrong{Traiter une couverture comme un simple coût. Elle est un coût ou un gain selon l'écart de taux, et le signe s'inverse pour un investisseur européen aux États-Unis.}
\end{solution}

\begin{solution}{iq:iv2:macro-market-knowledge-and-pitches:24}
Ils impliquent environ 64~\% d'une baisse lors de la première réunion, et environ 28~\% d'une baisse lors de la seconde.

\textbf{Première réunion.} La moyenne du mois est à moitié à 4.33 et à moitié au nouveau taux $R$~: $(4.33 + R)/2 = 4.25$. Donc $R = 4.17$, soit une baisse intégrée de 16 points de base, ou $16/25 = 64~\%$.

\textbf{Seconde réunion.} Le contrat à terme du mois 2 moyenne le taux après la seconde réunion~: 4.10. Cela fait 7 points de base de moins que 4.17, soit $7/25 = 28~\%$.

Au total, 23 points de base de baisses sont donc intégrés à l'horizon du mois 2.

\textbf{Ce qu'il faut ajuster.} Les décomptes de jours~: utiliser les jours exacts avant et après la réunion. La taille des baisses~: une baisse de 50 points de base peut aussi être intégrée. La prime de risque~: les contrats à terme ne sont pas une prévision pure. L'écart entre le taux effectif et le milieu de la fourchette cible. Et si le contrat se règle sur la moyenne du mois, comme c'est le cas ici.

\iqwrong{Comparer directement 4.25 à 4.33~: $8/25 = 32~\%$. Cela oublie que la réunion se tient au milieu du mois, de sorte que la moitié de la moyenne est déjà fixée.}
\end{solution}

\begin{solution}{iq:iv2:macro-market-knowledge-and-pitches:25}
Le P\&L vaut $+\$1.4$ million, $+\$0.3$ million et $-\$2$ millions. La moyenne est nulle, et l'écart type vaut environ 1.24 million de dollars. À 1~\%, le straddle ne vaut pas la peine d'être acheté sur ces chiffres~: son espérance de gain est d'environ 0.78~\%. Je réduirais la position en vendant des contrats à terme sur indice actions, ou en sortant du deux ans~: le livre ne constitue qu'un seul pari contre une inflation élevée.

\textbf{États.} Le P\&L en millions vaut $100r - 0.05\,\Delta y$~:
\begin{itemize}
\item 0.2~\%~: $1.0 + 0.4 = +1.4$.
\item 0.3~\%~: $0.3 + 0 = +0.3$.
\item 0.4~\%~: $-1.5 - 0.5 = -2.0$.
\end{itemize}

\textbf{Moyenne.} $0.25 \times 1.4 + 0.5 \times 0.3 + 0.25 \times (-2.0) = 0$.

\textbf{Dispersion.} La variance vaut $0.25 \times 1.96 + 0.5 \times 0.09 + 0.25 \times 4 = 1.535$, et sa racine vaut 1.239.

\textbf{Straddle.} L'amplitude absolue espérée du mouvement vaut $0.25 \times 1 + 0.5 \times 0.3 + 0.25 \times 1.5 = 0.775~\%$, pour un coût de 1~\%. Donc, selon les chiffres posés ici, on perd environ 22 cents par dollar de prime. Il ne rapporterait que si la carte des mouvements est trop étroite, ou si la distribution comporte plus de ces trois états.

\textbf{Lecture de la position.} Les actions et le deux ans perdent tous deux dans le mauvais état~: il n'y a donc aucune diversification, ils constituent le même pari. Le deux ans paraît petit à côté de 100 millions de dollars d'actions, mais il représente 25~\% du pire scénario.

\textbf{Ce soir.} La moyenne est nulle, et la queue gauche pèse davantage. Il n'y a donc aucune raison de garder cette taille pendant la publication. Il faut fixer sa perte maximale, disons 1~million de dollars, et raisonner à rebours. Réduire les actions de 30~\% ramène le mauvais état à environ $-1.55$ million de dollars. Clôturer aussi le deux ans le ramène à $-1.05$ million de dollars. Si le desk souhaite conserver la position longue, il faut acheter un put spread, et accepter un coût que l'on peut énoncer.

\textbf{Au responsable.} «~Nous sommes acheteurs d'un seul risque~: une inflation inférieure au consensus. Une mauvaise publication coûte environ 2 millions de dollars, avec une probabilité de 25~\% selon mes chiffres. Le reste de la distribution ne rémunère pas ce risque. Je propose de réduire d'environ un tiers, de clôturer le deux ans, et de rentrer à nouveau après la publication. Voici ce qui me ferait changer d'avis.~» Le découpage en états est un jugement~: il faut le dire.

\iqwrong{Comparer l'espérance de chaque position et s'arrêter là. Les deux jambes sont corrélées par l'événement, de sorte qu'une position qui paraît un mélange d'actions et d'obligations est en réalité un pari à un seul facteur.}
\iqfollow{La publication sort à 0.4~\% et les actions ne baissent que de 0.5~\%. Qu'en pensez-vous~?}{Le marché a moins bougé que la carte. Soit il était positionné pour cela, soit il lit autre chose dans la publication. Il ne faut pas supposer que le mouvement restant est à contrer. Il faut vérifier le deux ans~: s'il confirme la carte construite plus haut, la réaction des actions peut être en retard, mais il faut dire ce que l'on devrait voir.}
\end{solution}
